\documentclass[11pt,a4paper]{article}

\usepackage[utf8]{inputenc}
\usepackage[T1]{fontenc}
\usepackage{lmodern}
\usepackage{amsmath,amssymb,amsthm}
\usepackage{enumerate}
\usepackage[shortlabels]{enumitem}
\usepackage{hyperref}
\hypersetup{
  pdftitle={Shell-Level Saturation and the Conditional Fibre-to-Rate Prescription},
  pdfauthor={Jérôme Beau}
}
\usepackage{doi}

\input{glyphtounicode}
\pdfgentounicode=1

\newtheorem{theorem}{Theorem}[section]
\newtheorem{proposition}[theorem]{Proposition}
\newtheorem{corollary}[theorem]{Corollary}
\newtheorem{lemma}[theorem]{Lemma}
\theoremstyle{definition}
\newtheorem{definition}[theorem]{Definition}
\newtheorem{conjecture}[theorem]{Conjecture}
\newtheorem{openproblem}[theorem]{Open Problem}
\theoremstyle{remark}
\newtheorem{remark}[theorem]{Remark}

\newcommand{\bstar}{\beta^*}

\title{Shell-Level Saturation and the Conditional
  Fibre-to-Rate Prescription\\[0.5ex]
  \large O21 --- Spectral Admissibility Sub-programme}
\author{J.~Beau\\[0.5ex]
\small Independent Researcher}
\date{2026\\Version 1.2}

\begin{document}

  \maketitle

  \begin{abstract}
    The spectral admissibility sub-programme of Cosmochrony has
    identified the canonical fibre-level observable
    $\sigma_{\mathrm{pair}}^{\mathrm{can}}(n)$ as the relevant
    quantity for the lepton mass hierarchy.
    However, existing saturation criteria rely on the power-law
    fit $\sigma(n) \sim C\,n^{-\delta}$, whose amplitude $C$ and
    exponent $\delta$ depend on the fitting window and are not
    intrinsic.

    We construct a fully observable saturation rank
    $n_3^{\mathrm{obs}}(n_0, n_1, q)$ by eliminating both $C$
    and $\delta$ using two local measurements of
    $\sigma_{\mathrm{pair}}^{\mathrm{can}}$.
    This rank is invariant under global rescaling and all
    pipeline-dependent normalisations, and depends only on
    directly measurable quantities.

    We conjecture that physical admissibility requires
    $n_3^{\mathrm{obs}}$ to align with a BFS shell of the Cayley
    graph $\mathrm{Cay}(G_q, S_q)$, providing a geometric
    interpretation of saturation.
    The historical reciprocal image $1/(\delta_{\mathrm{pair}}+1/2)$ is retained as a
    defined numerical diagnostic only.
    O6 proves an upper growth bound and refutes equality under its own premises.
    Neither shell alignment nor a cross-substrate identification repairs that failed implication.
    The fibre construction does not supply a native Heisenberg
    growth equation carrying the pair observable.

    O22 and O24 discharge the shell-alignment and rank-three
    conditions on the fibre side of the chain.
    Those results do not supply the independent Heisenberg
    growth carrier required by the capacity-to-rate step.
    \textbf{Interpretive outlook.}
    A crossing-rank construction and a physical particle-hierarchy rate remain distinct; the
    reciprocal comparison supplies no rate-selection mechanism.
  \end{abstract}

  \paragraph{Keywords.}
    spectral admissibility; non-injective projection; fibre-level observables; Weil representation; Heisenberg graphs;
    projective capacity; power-law decay; intrinsic saturation criterion; BFS shell structure; scale invariance; lepton
    mass hierarchy; cascade exponent; subdiffusive growth

    \clearpage
    \tableofcontents

    \clearpage
  \section{Introduction}
  \label{sec:intro}

  The spectral admissibility sub-programme of the Cosmochrony
  framework~\cite{Cosmochrony} has progressively clarified a
  structural tension between the exact Weil-block capacity
  exponent and the phenomenological target set by the lepton
  mass hierarchy.
  Papers O13--O15~\cite{Beau2026a17,Beau2026a18,Beau2026a19}
  ruled out finite-size effects and aggregation corrections as
  resolutions; O16--O19~\cite{Beau2026a20,Beau2026a21,Beau2026a22,Beau2026a23}
  relocated the problem at the fibre level, establishing that
  the physically relevant observable is the canonical pair
  quantity $\sigma_{\mathrm{pair}}^{\mathrm{can}}(n)$
  constructed from conjugate Weil blocks $\{c, q-c\}$.
  Paper O20~\cite{Beau2026a24} proposed a projective persistence
  criterion for this observable, conditional on two structural
  hypotheses concerning its relationship to the Born--Infeld
  saturation envelope.

  A central difficulty that remains unresolved after O20 is
  the following.
  The power-law fit
  \[
    \sigma_{\mathrm{pair}}^{\mathrm{can}}(n)
    \approx C\,n^{-\delta_{\mathrm{pair}}}
  \]
  introduces two fitted parameters, the amplitude $C$ and the
  exponent $\delta_{\mathrm{pair}}$, both of which depend on
  the fitting window and therefore on pipeline conventions.
  Any admissibility criterion formulated in terms of these
  parameters inherits this dependence and is therefore not
  intrinsic.

  The central problem addressed in this paper is therefore:
  can the saturation condition of O20 be reformulated in a
  way that is entirely independent of fitted parameters and
  normalisation choices?

  We construct a positive answer under the power-law
  hypothesis.
  By combining two local observations of
  $\sigma_{\mathrm{pair}}^{\mathrm{can}}$ at reference shells
  $n_0$ and $n_1$, one can eliminate both $C$ and
  $\delta_{\mathrm{pair}}$ simultaneously, producing a fully
  observable saturation rank $n_3^{\mathrm{obs}}(n_0, n_1, q)$
  that is invariant under all pipeline-dependent
  transformations identified in
  O17--O19~\cite{Beau2026a21,Beau2026a22,Beau2026a23}
  and, in particular, under the amplitude artefact
  $r(c,q) \in \{1,2,3,4\}$.

  The key conceptual step is then to interpret saturation not
  as the attainment of a numerical threshold in a continuous
  fit, but as a \emph{geometric alignment condition} with the
  discrete shell structure of the Cayley graph
  $\mathrm{Cay}(G_q, S_q)$.
  This leads to the shell-alignment conjecture: a conjugate
  pair is physically admissible if and only if
  $n_3^{\mathrm{obs}}$ aligns with a BFS shell
  (Conjecture~\ref{conj:shell-saturation}).

  This reformulation has three structural consequences.
  First, it removes all free parameters from the admissibility
  criterion, making it directly testable across primes and
  pairs without pipeline recalibration.
  Second, it reveals a resonance condition between two a priori
  independent structures: the continuous power-law decay
  induced by Weil representation dynamics and the discrete
  shell geometry of the Heisenberg Cayley graph.
  Third, it provides a structural interpretation of the
  threshold $9 = 3^2$ as the normalised product of the
  fibre-level cumulative spans at the three-dimensional
  saturation point, connecting the admissibility condition to
  the emergence of three physically stable directions under
  the non-injective projection~$\Pi$.

  The criterion constructed here is intrinsic, observable, and
  independent of representation choices.
  It provides an intrinsic fibre-side ingredient for the
  candidate $\delta_{\mathrm{pair}} \to \beta^*$ pathway.
  It does not derive the growth-equation side of that pathway
  on the Heisenberg substrate~\cite{Beau2026sgn}.

  The paper is organised as follows.
  Section~\ref{sec:background} recalls the relevant results of
  O16--O20.
  Section~\ref{sec:observable-criterion} constructs
  $n_3^{\mathrm{obs}}$ and establishes its invariance
  properties.
  Section~\ref{sec:shell-conjecture} states the shell-alignment
  conjecture and discusses its numerical status.
  Section~\ref{sec:transfer} withdraws the former transfer theorem and defines the reciprocal
  image solely as a legacy numerical diagnostic.
  Section~\ref{sec:conclusion} discusses open directions.

  \section{Background: From Block to Fibre Level}
  \label{sec:background}

  \subsection{Block-level observable and its limitations}

    Papers O12 and O13~\cite{Beau2026a16,Beau2026a17} introduced
    the Gram--Schmidt span tracker on Heisenberg Cayley graphs
    $\mathrm{Cay}(G_q, S_q)$ with
    $G_q = \mathrm{Heis}_3(\mathbb{Z}/q\mathbb{Z})$ for prime
    $q$, establishing the mean block-level observable
    $\bar{\sigma}(n)$ with decay exponent
    $\hat{\delta}_{\mathrm{exact}} \approx 3.72$.
    O14 and O15~\cite{Beau2026a18,Beau2026a19} demonstrated that
    this exponent lies structurally below the phenomenological
    target $\delta \in [7.4, 10.6]$, and that no aggregation
    of block data can bridge this gap (O15 Corollary~4.8).

  \subsection{Fibre-level resolution}

    O16~\cite{Beau2026a20} identified the resolution: the
    physically relevant observable is defined on conjugate pairs
    $\{c, q-c\}$, giving
    $\sigma_{\mathrm{pair}}(n) = \sigma_c(n)\,\sigma_{q-c}(n)$
    with exponent $\delta_{\mathrm{pair}} = 2\delta_c \approx
  7.44$.
    O17~\cite{Beau2026a21} established the structural derivation:
    $\delta_{q-c} = \delta_c$ follows from
    $\rho_{q-c} = \overline{\rho_c}$ and norm-invariance of
    Gram--Schmidt.
    O18~\cite{Beau2026a22} establishes the parity of the Born--Infeld
    action as a covariance of the response family and states the identification of
    $\{c, q-c\}$ with a fibre of $\Pi$ as an open problem; the pairing
    keeps the status of a structurally motivated hypothesis.
    O19~\cite{Beau2026a23} constructed the canonical observable
    $\sigma_{\mathrm{pair}}^{\mathrm{can}}(n)$, eliminating the
    residual amplitude factor $r(c,q) \in \{1,2,3,4\}$.

    \begin{remark}[Observable was mis-identified, not mis-measured]
      \label{rem:mis-identified-background}
      The move from block to fibre level is not a correction of
      O12--O15 but a change of observable class.
      O15 closes the block-level programme; O16--O19 open and
      populate the fibre-level programme.
    \end{remark}

  \subsection{The open problem after O20}

    O20~\cite{Beau2026a24} proposed a projective persistence
    criterion for $\sigma_{\mathrm{pair}}^{\mathrm{can}}$,
    conditional on two hypotheses:
    that the observable is a monotone function of an effective
    fibre-level amplitude~[H1], and that this amplitude inherits
    the Born--Infeld saturation bound~[H2].
    The central open problem identified in O20 is the intrinsic
    derivation of the saturation threshold $\sigma_{\mathrm{BI}}$
    from the Born--Infeld constraint, without reference to the
    phenomenological window $\beta^* \in (0.09, 0.13)$.

    The present paper addresses this problem from a different
    angle: rather than deriving $\sigma_{\mathrm{BI}}$
    dynamically, we construct a saturation rank
    $n_3^{\mathrm{obs}}$ that is free of fitted parameters and
    interpret its geometric meaning.

  \section{A Fully Observable Saturation Criterion}
  \label{sec:observable-criterion}

  \subsection{Eliminating the amplitude and the exponent}

    The power-law fit
    $\sigma_{\mathrm{pair}}^{\mathrm{can}}(n) \approx
    C\,n^{-\delta_{\mathrm{pair}}}$ introduces two fitted
    parameters, $C$ and $\delta_{\mathrm{pair}}$, both of which
    depend on the fitting window and therefore on pipeline
    conventions.
    We eliminate both in succession.

    \paragraph{Eliminating $C$.}
      For any two shells $n_0 < n_3$ in the power-law regime,
      \[
        \frac{\sigma(n_3)}{\sigma(n_0)}
        = \left(\frac{n_3}{n_0}\right)^{-\delta}.
      \]
      Setting $\sigma(n_3) = 9/q^2$ and solving for $n_3$:
      \begin{equation}
        \label{eq:n3-no-C}
        n_3
        = n_0 \left(\frac{q^2\,\sigma(n_0)}{9}\right)^{1/\delta}.
      \end{equation}
      The amplitude $C$ has cancelled.

    \paragraph{Eliminating $\delta$.}
      The local exponent between any two reference shells
      $n_0 < n_1$ within the power-law regime is
      \begin{equation}
        \label{eq:delta-local}
        \delta(n_0, n_1)
        = \frac{\log\bigl(\sigma(n_0)/\sigma(n_1)\bigr)}
        {\log(n_1/n_0)}.
      \end{equation}
      Substituting~\eqref{eq:delta-local}
      into~\eqref{eq:n3-no-C} gives the fully observable
      expression:
      \begin{equation}
        \label{eq:n3-observable}
        \boxed{
          n_3^{\mathrm{obs}}(n_0, n_1, q)
          \;=\;
          n_0
          \left(\frac{q^2\,\sigma(n_0)}{9}\right)^{%
            \frac{\log(n_1/n_0)}{\log(\sigma(n_0)/\sigma(n_1))}}
        }
      \end{equation}
      This expression contains no fitted parameters.
      It depends only on the observable values
      $\sigma_{\mathrm{pair}}^{\mathrm{can}}(n_0)$ and
      $\sigma_{\mathrm{pair}}^{\mathrm{can}}(n_1)$ at two reference
      shells, the prime $q$, and the integer $9$.

      \begin{definition}[Observable saturation rank]
        \label{def:n3obs}
        Let $n_0 < n_1$ be two shells in the power-law regime of
        $\sigma_{\mathrm{pair}}^{\mathrm{can}}$.
        The \emph{observable saturation rank} is
        \[
          n_3^{\mathrm{obs}}(n_0, n_1, q)
          \;:=\;
          n_0
          \left(\frac{q^2\,\sigma_{\mathrm{pair}}^{\mathrm{can}}(n_0)}{9}
          \right)^{%
            \frac{\log(n_1/n_0)}{\log\!\left(
                                         \sigma_{\mathrm{pair}}^{\mathrm{can}}(n_0)\,/\,
                                         \sigma_{\mathrm{pair}}^{\mathrm{can}}(n_1)
            \right)}
          }.
        \]
      \end{definition}

      \begin{proposition}[Local power-law consistency]
        \label{prop:local-consistency}
        If $\sigma_{\mathrm{pair}}^{\mathrm{can}}(n)$ follows an
        exact power law on an interval $[n_0, n_1]$, then
        $n_3^{\mathrm{obs}}(n_0, n_1, q)$ is independent of the
        choice of $(n_0, n_1)$ within that interval.
        Conversely, any observed dependence of $n_3^{\mathrm{obs}}$
        on $(n_0, n_1)$ directly measures deviations from an exact
        power-law regime.
      \end{proposition}

      \begin{proof}
        Under the exact power law $\sigma(n) = C n^{-\delta}$,
        the local exponent $\delta(n_0, n_1) = \delta$ for any
        $n_0 < n_1$ in the regime.
        Substituting into Definition~\ref{def:n3obs}:
        \[
          n_3^{\mathrm{obs}}
          = n_0 \left(\frac{q^2 C n_0^{-\delta}}{9}\right)^{1/\delta}
          = \left(\frac{q^2 C}{9}\right)^{1/\delta},
        \]
        which is independent of $n_0$ and $n_1$.
        The converse follows immediately.
      \end{proof}

      \begin{remark}[Invariance under global rescaling]
        \label{rem:invariance}
        $n_3^{\mathrm{obs}}$ is invariant under $\sigma \mapsto
        \lambda \sigma$ for any constant $\lambda > 0$: the
        exponent in Definition~\ref{def:n3obs} uses only the ratio
        $\sigma(n_0)/\sigma(n_1)$, which is unaffected by a global
        multiplicative factor.
        In particular, $n_3^{\mathrm{obs}}$ is independent of the
        pipeline normalisation factor $r(c,q) \in \{1,2,3,4\}$
        identified in
        O17--O19~\cite{Beau2026a21,Beau2026a22,Beau2026a23}.
        The criterion therefore depends only on the \emph{shape} of
        the decay, not on its absolute normalisation.
      \end{remark}

  \section{The Shell-Alignment Conjecture}
  \label{sec:shell-conjecture}

  Here $S_n = \{g \in G_q \mid d(e, g) = n\}$ denotes the
  sphere of radius $n$ in $\mathrm{Cay}(G_q, S_q)$, where $d$
  is the word metric.

  \begin{conjecture}[Shell-level saturation]
    \label{conj:shell-saturation}
    A conjugate pair $\{c, q-c\}$ is physically admissible if
    and only if the observable saturation rank aligns with a
    shell of the Cayley graph $\mathrm{Cay}(G_q, S_q)$:
    \[
      \mathrm{dist}\!\left(
                       n_3^{\mathrm{obs}}(n_0, n_1, q),\;
                       \{n \mid S_n \neq \varnothing\}
      \right)
      \;\ll\; 1,
    \]
    where $S_n = \{g \in G_q \mid d(e, g) = n\}$ denotes the
    BFS shell at depth $n$ in $\mathrm{Cay}(G_q, S_q)$, and $d$
    is the word metric.
  \end{conjecture}

  \begin{remark}[Physical interpretation and logarithmic
  commensurability]
    \label{rem:physical-interp}
    The Cayley graph $\mathrm{Cay}(G_q, S_q)$ imposes a discrete
    shell structure $S_0, S_1, S_2, \ldots$ determined by the
    group geometry of $\mathrm{Heis}_3(\mathbb{Z}/q\mathbb{Z})$.
    The power-law decay of $\sigma_{\mathrm{pair}}^{\mathrm{can}}$
    is governed by the Weil representation dynamics.
    These two structures are a priori independent.
    Conjecture~\ref{conj:shell-saturation} asserts that physical
    admissibility requires their \emph{resonance}: the point at
    which the capacity observable reaches the saturation
    threshold must coincide with a shell boundary.
    This is a condition on the representation theory of
    $\mathrm{Heis}_3(\mathbb{Z}/q\mathbb{Z})$, not a consequence
    of any choice of numeration.

    Writing $n_3^{\mathrm{obs}} = n_0 \cdot A^{1/\delta}$ where
    $A = q^2\,\sigma(n_0)/9 > 0$, the condition can be viewed
    as a \emph{logarithmic commensurability constraint}: the
    local decay rate $\delta$ and the amplitude ratio $A$ must
    be commensurate with the integer shell structure of the
    graph.
  \end{remark}

  \begin{remark}[The origin of the threshold $9 = 3^2$]
    \label{rem:origin-9}
    The number $9 = 3^2$ in Definition~\ref{def:n3obs} is not
    an external input.
    If $\Sigma_c(n_3) = \Sigma_{q-c}(n_3) = 3$ at the
    saturation shell, then
    \[
      q^2\,\sigma_{\mathrm{pair}}^{\mathrm{can}}(n_3)
      \;=\;
      \frac{\Sigma_c(n_3)}{q} \cdot
      \frac{\Sigma_{q-c}(n_3)}{q} \cdot q^2
      \;=\; 3^2 \;=\; 9.
    \]
    The threshold $9/q^2$ is therefore the normalised product of
    the two fibre-level cumulative spans at the
    three-dimensional saturation point.
    This identification is independent of any basis choice,
    as it involves only cumulative span dimensions.
    The integer $3$ reflects the number of physically stable
    directions selected by the projection $\Pi$ in the Weil
    realisation of $\mathrm{Heis}_3(\mathbb{Z}/q\mathbb{Z})$;
    whether $\Sigma_c(n_3) = 3$ is a derived consequence of the
    Born--Infeld structure or an independent arithmetic property
    of the Weil representation is an open problem.
  \end{remark}

  \begin{remark}[Stability as a diagnostic]
    \label{rem:stability-diagnostic}
    By Proposition~\ref{prop:local-consistency}, the stability
    of $n_3^{\mathrm{obs}}$ under variation of $(n_0, n_1)$
    within the pre-saturation window provides a direct empirical
    test of the power-law hypothesis.
    A stable $n_3^{\mathrm{obs}}$ confirms the regime;
    instability quantifies finite-size corrections and guides
    the choice of fitting window for larger~$q$.
  \end{remark}

  \subsection{Numerical status}

    The shell-alignment conjecture makes a concrete and
    falsifiable prediction: for physically admissible conjugate
    pairs, $n_3^{\mathrm{obs}}(n_0, n_1, q)$ should lie close
    to an admissible BFS shell depth.

    \begin{table}[ht]
      \centering
      \begin{tabular}{cccccc}
        \hline
        $q$ & $(c,\,q-c)$ & $\delta_{\mathrm{pair}}$ & $R^2$ &
        $n_3^{\mathrm{obs}}$ &
        $\mathrm{dist}(n_3^{\mathrm{obs}}, \mathbb{Z})$ \\
        \hline
        61 & $(1,\,60)$ & 7.512 & 1.0000 & 4.80 & 0.20 \\
        \hline
      \end{tabular}
      \caption{Observable saturation rank for available pairs with
        $\delta_{\mathrm{pair}} \in [7.4, 10.6]$.
        $n_3^{\mathrm{obs}}$ is computed from
        Definition~\ref{def:n3obs} with $(n_0, n_1) = (2, 7)$,
        the pre-saturation window of~\cite{Beau2026a17}.
        The distance to the nearest BFS shell depth provides a
        direct test of Conjecture~\ref{conj:shell-saturation}.}
      \label{tab:n3obs}
    \end{table}

    At integer shells, the threshold $9/q^2$ is bracketed as:
    \[
      q^2\,\sigma_{\mathrm{pair}}^{\mathrm{can}}(4) \approx 6.64,
      \qquad
      q^2\,\sigma_{\mathrm{pair}}^{\mathrm{can}}(5) \approx 1.66,
    \]
    showing that $n_3^{\mathrm{obs}} \approx 4.80$ lies strictly
    between two consecutive shells.

    The value $\mathrm{dist}(n_3^{\mathrm{obs}}, \mathbb{Z})
  \approx 0.20$ is intermediate: it does not indicate exact
    alignment, but is also significantly below maximal
    misalignment ($0.5$).

    Two scenarios are distinguishable by further data.
    The conjecture predicts not that any individual pair has
    $\mathrm{dist}(n_3^{\mathrm{obs}}, \mathbb{Z}) = 0$
    exactly, but that the \emph{distribution} of distances
    across physically admissible pairs concentrates near $0$,
    in contrast to a uniform distribution on $[0, \tfrac{1}{2}]$
    expected under the null hypothesis.
    In the \emph{alignment scenario}, this concentration becomes
    statistically significant as the dataset grows in $q$ and in
    the number of pairs.
    In the \emph{non-alignment scenario}, the distribution
    remains uniform, falsifying the conjecture.
    The single available data point
    ($\mathrm{dist} \approx 0.20$ for $q=61$) is consistent with
    both scenarios and does not discriminate between them.

    A decisive test requires computing
    $n_3^{\mathrm{obs}}$ for all conjugate pairs with
    $\delta_{\mathrm{pair}}$ in the historical range $[7.4,10.6]$ across
    $q \in \{101, 151, 211\}$, using the canonical pair
    observable of O19~\cite{Beau2026a23}.

  \section{The Legacy $\delta_{\mathrm{pair}} \to \beta^*$
  Diagnostic and the Withdrawn Transfer Claim}
  \label{sec:transfer}

  \subsection{Status of the O7 reciprocal comparison}

    The former O7 equality is replaced here by a definition of a legacy diagnostic:
    \begin{equation}
      \label{eq:O7-relation}
      \beta_{\mathrm{diag}}(\delta):=\frac{1}{\delta+\tfrac12}.
    \end{equation}
    O6~\cite{Beau2026a10} proves only $\beta_{\mathrm{eff}}\leq1/(\alpha+1/2)$ under its
    effective growth hypothesis (G$_R$) and $\mathcal R_n\sim p(n)^{-\alpha}$.
    Its exact power-law countermodel $\alpha=8$, $p(n)=n^{1/17}$ satisfies those premises
    while $\beta=1/17<2/17$ lies outside $(0.09,0.13)$.
    Thus equality fails before any transport to a Heisenberg pair observable is considered.
    A decay fit, a wide window, a determinate crossing rank and a cross-substrate identification
    do not establish the missing lower growth bound or select equality.
    The Span-Growth Note supplies the separate native-transfer obstruction for the corpus-defined
    Heisenberg objects~\cite{Beau2026sgn}.

  \subsection{Withdrawal of the transfer theorem}

    \begin{remark}[Former transfer claim withdrawn]
      \label{thm:conditional-transfer}
      The former theorem asserting a reciprocal cascade exponent under shell alignment and an
      external identification is withdrawn.
      Its conclusion was not a consequence of the cited O6/O7 premises.
      For a supplied pair exponent one may define
      \begin{equation}
        \label{eq:transfer}
        \beta^*_{\mathrm{presc}}:=\frac{1}{\delta_{\mathrm{pair}}+\tfrac12}.
      \end{equation}
      This is a legacy comparison value, not the physical cascade exponent $\beta^*$ and not
      a theorem conditional on the existence of a transport map.
    \end{remark}

    \begin{remark}[Independent status of the rank construction]
      \label{rem:conditional}
      This withdrawal does not change the two-point crossing formula, the stated shell-alignment
      conjecture or the reported finite-prime observations.
      Those objects supply no reciprocal rate equality.
      No additional saturation hypothesis is introduced to restore the withdrawn theorem.
    \end{remark}

  \subsection{Numerical evaluation}

    Substituting $\delta_{\mathrm{pair}} \approx 7.44$
    into the defined legacy diagnostic gives
    \begin{equation}
      \label{eq:beta-numerical}
      \beta^*_{\mathrm{presc}}
      := \frac{1}{7.44 + 0.5}
      \approx 0.126,
    \end{equation}
    which lies near the upper boundary of the phenomenological
    window $\beta^* \in (0.09, 0.13)$ established
    in~O3~\cite{Beau2026a7}.

    \begin{remark}[Position within the window]
      \label{rem:window-position}
      The value $\beta^*_{\mathrm{presc}} \approx 0.126$ is a
      phenomenological coincidence check near the upper boundary
      of the target, not a structural derivation.
      The interval $[7.4,10.6]$ is a historical comparison range, not a selected physical sub-spectrum.
      A larger supplied $\delta_{\mathrm{pair}}$ decreases the diagnostic reciprocal; it does
      not determine the physical $\beta^*$ or guarantee membership in its open window.
      The value $7.44$ uses the production convention
      $\log(n+1)$.
      On the same finite windows, translation to the paper-level
      $\log n$ definition gives approximately $6.6$--$6.9$ and
      therefore falls below $7.4$; this estimator conversion is
      conditional on those windows~\cite{Beau2026sgn}.
    \end{remark}

  \subsection{The role of the fibre doubling}

    The exact doubling $\delta_{\mathrm{pair}} = 2\delta_c$
    from O16--O17~\cite{Beau2026a20,Beau2026a21} gives the
    equivalent block-level expression:
    \begin{equation}
      \label{eq:transfer-block}
      \beta^*_{\mathrm{presc}}
      = \frac{1}{2\delta_c + \tfrac{1}{2}}.
    \end{equation}
    The factor $2$ in the pair exponent is structural rather than
    combinatorial: it reflects the independent and equal
    contribution of the two conjugate blocks to the pair-level
    capacity, established by $\delta_{q-c} = \delta_c$
    (O17 Theorem~2.1~\cite{Beau2026a21}, recalled in O18
    Theorem~3.1~\cite{Beau2026a22}).
    This exact doubling does not by itself justify inserting
    $\delta_{\mathrm{pair}}$ into the O7 growth equation.

  \subsection{Resolution of the O15 tension}

    O15 showed that the transfer $\delta \to \beta^*$ fails at
    the block level: no aggregation of block data recovers the
    fibre-level exponent.
    The present paper resolves the amplitude-dependent
    saturation criterion by defining $n_3^{\mathrm{obs}}$
    intrinsically at the fibre level.
    It does not resolve the separate growth-equation obstruction:
    no native process carrying
    $\sigma_{\mathrm{pair}}^{\mathrm{can}}$ is derived.

  \section{Conclusion}
  \label{sec:conclusion}

  This paper establishes a fully observable and intrinsic
  criterion for fibre-level admissibility in the spectral
  programme of Cosmochrony.
  By eliminating both the amplitude and the exponent from the
  power-law description of
  $\sigma_{\mathrm{pair}}^{\mathrm{can}}(n)$, we constructed
  the saturation rank $n_3^{\mathrm{obs}}$, which depends only
  on directly measurable values of the observable and is
  invariant under all pipeline-dependent normalisations.
  Under the shell-alignment conjecture, this rank acquires a
  discrete geometric interpretation as a BFS shell depth in
  $\mathrm{Cay}(G_q, S_q)$, providing a well-defined and
  scale-independent saturation event.
  The former transfer theorem is withdrawn: the additional cross-substrate identification
  does not repair the reciprocal equality refuted by O6's countermodel.
  The defined legacy image $\beta^*_{\mathrm{presc}}\approx0.126$ is a numerical comparison
  for a supplied exponent, not a consequence of shell alignment or a determination of $\beta^*$.
  The native Heisenberg transfer obstruction remains separate~\cite{Beau2026sgn}.

  The construction resolves the intrinsic-definition problem
  for the fibre saturation rank.
  It leaves the growth-equation transfer unresolved and, for the
  corpus-defined Heisenberg objects, refuted as a native
  composite.

  O22 proves shell alignment, and O24 closes the rank-three
  condition with respect to the fibre structure
  alone~\cite{Beau2026a26,Beau2026a28}.
  These fibre-side statements do not select a rate or restore the withdrawn reciprocal theorem.
  A new growth carrier and an independently justified rate-selection mechanism would be needed;
  neither follows from the crossing rank or from a transport identification alone.


  \bibliographystyle{plain}
      \bibliography{cosmochrony-bibliography}

\end{document}
